\begin{lemma}[Lemma 2.1]
Let $(y_h)_{h=0,1,2,\dots}$ be a sequence of nonnegative real numbers such that
for some positive numbers $c$, $\varepsilon$ and some $b>1$ the recursion
relation
\[
y_{h+1}\le c\,b^{h}\,y_h^{1+\varepsilon},\qquad h=0,1,2,\dots
\]
is satisfied.  Then, if
\[
y_0\le\theta=c^{-\frac1\varepsilon}\,b^{-\frac1{\varepsilon^2}},
\]
it follows that $\lim_{h\to\infty}y_h=0$.
\end{lemma}
\begin{proof}
This is Lemma~Lemma~2.1 of the thesis of the paper with the correction prescribed by the
paper's Note for the formalization: the recursion is $y_{h+1}\le
c\,b^{h}y_h^{1+\varepsilon}$, as in [LU, p.~66] and as the recursion is actually
applied in (2.6) of the thesis, rather than the printed $y_{h+1}\le
c\,b^{h+1}y_h^{1+\varepsilon}$, for which the threshold $\theta$ is
insufficient.  All powers are real powers of positive reals.

Set $\theta=c^{-1/\varepsilon}b^{-1/\varepsilon^2}$; since $c>0$ and $b>0$ we
have $\theta>0$, and in particular $\theta\ge0$.

\emph{Step 1 (the threshold identity).}  We claim
\[
c\,\theta^{\varepsilon}=b^{-1/\varepsilon},
\]
with exact equality.  Indeed, since $c^{-1/\varepsilon}\ge0$ and
$b^{-1/\varepsilon^2}\ge0$,
\[
\theta^{\varepsilon}
=\bigl(c^{-1/\varepsilon}\bigr)^{\varepsilon}\bigl(b^{-1/\varepsilon^{2}}\bigr)^{\varepsilon}
=c^{(-1/\varepsilon)\cdot\varepsilon}\,b^{(-1/\varepsilon^{2})\cdot\varepsilon}
=c^{-1}\,b^{-1/\varepsilon},
\]
using that $(xy)^{\varepsilon}=x^{\varepsilon}y^{\varepsilon}$ for
$x,y\ge0$ and $(x^{s})^{t}=x^{st}$ for $x\ge0$, together with
$(-1/\varepsilon)\varepsilon=-1$ and
$(-1/\varepsilon^{2})\varepsilon=-1/\varepsilon$ (here $\varepsilon\neq0$).
Multiplying by $c$ and using $c\cdot c^{-1}=1$ (valid as $c\neq0$) gives the
claim.

\emph{Step 2 (the induction invariant).}  We prove by induction on $h$ that
\[
y_h\le\theta\,b^{-h/\varepsilon}\qquad\text{for all }h=0,1,2,\dots
\]
For $h=0$ this is the hypothesis $y_0\le\theta$, because
$b^{-0/\varepsilon}=b^{0}=1$.

Assume $y_k\le\theta b^{-k/\varepsilon}$ for some $k\ge0$.  Both sides are
nonnegative ($y_k\ge0$ by hypothesis, and $\theta b^{-k/\varepsilon}\ge0$), and
the exponent $1+\varepsilon$ is positive, so raising to the power
$1+\varepsilon$ preserves the inequality:
\[
y_k^{1+\varepsilon}\le\bigl(\theta\,b^{-k/\varepsilon}\bigr)^{1+\varepsilon}.
\]
Since $c\,b^{k}\ge0$, multiplying by $c\,b^{k}$ and combining with the recursion
at $h=k$ gives
\[
y_{k+1}\le c\,b^{k}\,y_k^{1+\varepsilon}
\le c\,b^{k}\bigl(\theta\,b^{-k/\varepsilon}\bigr)^{1+\varepsilon}.
\]
We now compute the right-hand side exactly.  First, since
$\theta\ge0$ and $b^{-k/\varepsilon}\ge0$,
\[
\bigl(\theta\,b^{-k/\varepsilon}\bigr)^{1+\varepsilon}
=\theta^{1+\varepsilon}\,b^{(-k/\varepsilon)(1+\varepsilon)}
=\theta\cdot\theta^{\varepsilon}\cdot b^{(-k/\varepsilon)(1+\varepsilon)} ,
\]
using $\theta^{1+\varepsilon}=\theta^{1}\theta^{\varepsilon}=\theta\,\theta^{\varepsilon}$.
Therefore
\[
c\,b^{k}\bigl(\theta\,b^{-k/\varepsilon}\bigr)^{1+\varepsilon}
=\theta\cdot\bigl(c\,\theta^{\varepsilon}\bigr)\cdot
  \Bigl(b^{k}\,b^{(-k/\varepsilon)(1+\varepsilon)}\Bigr)
=\theta\cdot b^{-1/\varepsilon}\cdot b^{\,k+(-k/\varepsilon)(1+\varepsilon)} ,
\]
where we used Step 1 and $b^{s}b^{t}=b^{s+t}$ (valid since $b>0$).  The
remaining exponent simplifies: since $\varepsilon\neq0$,
\[
k+\Bigl(-\frac k\varepsilon\Bigr)(1+\varepsilon)
=k-\frac k\varepsilon-k=-\frac k\varepsilon .
\]
Hence
\[
c\,b^{k}\bigl(\theta\,b^{-k/\varepsilon}\bigr)^{1+\varepsilon}
=\theta\,b^{-1/\varepsilon}\,b^{-k/\varepsilon}
=\theta\,b^{-(k+1)/\varepsilon},
\]
again by $b^{s}b^{t}=b^{s+t}$ and
$-1/\varepsilon-k/\varepsilon=-(k+1)/\varepsilon$.  Combining,
$y_{k+1}\le\theta\,b^{-(k+1)/\varepsilon}$, which completes the induction.

Observe that the chain of identities above is a chain of \emph{equalities}: the
induction closes with no margin to spare.  This is exactly why the printed form
of the recursion, with $b^{h+1}$ in place of $b^{h}$, fails with this
$\theta$: it would produce an extra factor $b>1$ in the last display and break
the induction.

\emph{Step 3 (passage to the limit).}  Put $q=b^{-1/\varepsilon}$.  Then $q>0$,
and since $b>1$ and $-1/\varepsilon<0$ we have $q<1$.  Moreover, for every
$h\in\mathbb N$,
\[
b^{-h/\varepsilon}=\bigl(b^{-1/\varepsilon}\bigr)^{h}=q^{h},
\]
because $b>0$ and $(-1/\varepsilon)\cdot h=-h/\varepsilon$.  Thus Step 2 reads
\[
0\le y_h\le\theta\,q^{h}\qquad\text{for all }h,
\]
the left inequality being the hypothesis that the $y_h$ are nonnegative.  Since
$0\le q<1$, the geometric sequence $q^h$ tends to $0$, hence so does
$\theta q^{h}$.  By the squeeze theorem, $y_h\to0$ as $h\to\infty$.
\end{proof}
